\subsection{Relations and Equivalence}

\begin{definition} \label{def:relations} \hlt{Relations}\\
Let $X,Y$ be sets.
\begin{enumerate}[label=(\roman*)]
\setlength{\itemsep}{0pt}
\item \hlt{From $X$ to $Y$}: A subset $R\subseteq X\times Y$ is a relation from $X$ to $Y$.
\item \hlt{Relation notation}: Given a relation $R$ from $X$ to $Y$, write $xRy\Longleftrightarrow(x,y)\in R$ for $(x,y)\in X\times Y$.
\item \hlt{On $X$}: A subset $R\subseteq X^2$ is a relation on $X$.
\end{enumerate}
\end{definition}

\begin{definition} \label{def:equivalence-relation} \hlt{Equivalence Relations}\\
Let $X$ be a set and let $\sim$ be a relation on $X$.
\begin{enumerate}[label=(\roman*)]
\setlength{\itemsep}{0pt}
\item \hlt{Reflexive}: $\sim$ is reflexive iff $(\forall x\in X)\,x\sim x$.
\item \hlt{Symmetric}: $\sim$ is symmetric iff $(\forall x,y\in X)\,[x\sim y\Rightarrow y\sim x]$.
\item \hlt{Transitive}: $\sim$ is transitive iff $(\forall x,y,z\in X)\,[(x\sim y\land y\sim z)\Rightarrow x\sim z]$.
\item \hlt{Equivalence relation}: $\sim$ is an equivalence relation iff conditions (i)--(iii) hold.
\end{enumerate}
\end{definition}

\begin{definition} \label{def:equivalence-classes-invariants-canonical-forms} \hlt{Equivalence Classes, Invariants, and Canonical Forms}\\
Let $X,Y$ be sets and let $\sim$ be an equivalence relation on $X$.
\begin{enumerate}[label=(\roman*)]
\setlength{\itemsep}{0pt}
\item \hlt{Equivalence class}: For $x\in X$, set $[x]_{\sim}\defeq\{y\in X:y\sim x\}$.
\item \hlt{Invariant}: A map $f:X\to Y$ is an invariant of $\sim$ iff $(\forall x,y\in X)\,[x\sim y\Rightarrow f(x)=f(y)]$.
\item \hlt{Complete invariant}: A map $f:X\to Y$ is a complete invariant of $\sim$ iff $(\forall x,y\in X)\,[x\sim y\Longleftrightarrow f(x)=f(y)]$.
\item \hlt{Canonical forms}: A subset $\mathcal{C}\subseteq X$ is a set of canonical forms for $\sim$ iff $(\forall x\in X)(\exists! c\in\mathcal{C})\,c\sim x$.
\end{enumerate}
\end{definition}
